% !TeX root=../main.tex

\university{دانشگاه تهران}
\college{پردیس دانشکده‌های فنی}
\faculty{دانشکدهٔ مهندسی برق و کامپیوتر}
\department{گروه مهندسی کامپیوتر}
\subject{مهندسی کامپیوتر}
\field{نرم‌افزار}

\title{کاهش مدل معنایی ربکا با استفاده از رابطهٔ شبیه‌سازی دوسویهٔ ضعیف زمان‌دار}

\firstsupervisor{دکتر فاطمه قاسمی اصفهانی}
\firstsupervisorrank{استادیار}
\internaljudge{دکتر رامتین خسروی}
\internaljudgerank{}

\name{امیرپارسا}
\surname{موبد}
\studentID{۸۱۰۱۰۰۲۷۱}
\thesisdate{شهریور ۱۴۰۵}

\faabstract{
وارسی مدل با پیمایش فضای حالت یک سامانه، ویژگی‌های مورد انتظار را به‌صورت خودکار بررسی می‌کند. در مدل‌های هم‌روند و زمان‌دار، ترکیب ترتیب اجرای مؤلفه‌ها، صف پیام‌ها و محدودیت‌های زمانی می‌تواند به رشد سریع فضای حالت و افزایش هزینهٔ تحلیل بینجامد. کاهش مدل، در صورتی که رفتار مورد بررسی را حفظ کند، راهی برای کنترل این مسئله است.

در این پروژه، کاهش مدل معنایی زمان‌دار ربکا بر اساس شبیه‌سازی دوسویهٔ ضعیف زمان‌دار بررسی و پیاده‌سازی شده است. در این روش، تعامل‌های نامرتبط با سطح مشاهدهٔ انتخاب‌شده پنهان می‌شوند و حالت‌هایی که از نظر رفتار قابل مشاهده و گذر زمان تمایزپذیر نیستند در یک رده قرار می‌گیرند. سپس از این رده‌ها مدلی کوچک‌تر ساخته می‌شود. ارزیابی روی نمونه‌های موجود نشان داد که این روش می‌تواند ضمن حفظ معیار رفتاری انتخاب‌شده، اندازهٔ مدل را کاهش دهد. نتیجهٔ پروژه مبنایی برای افزودن این نوع کاهش به فرایند تحلیل مدل‌های ربکا در افرا فراهم می‌کند.
}

\keywords{ربکا، افرا، شبیه‌سازی دوسویهٔ ضعیف زمان‌دار، روش‌های صوری، وارسی مدل}
